\chapter{Il modulo \code{concurrency-control}}
\label{cap:12-concurrency-control}

Questo \`e il modulo su cui il progetto ha investito di pi\`u, ed \`e il capitolo
con il maggior contenuto originale del documento. Il modulo decide, per ciascuna
funzione e a intervalli regolari, quante invocazioni possono essere
contemporaneamente in volo. \`E 1\,314 righe in quindici classi --- pi\`u di un
quarto del codice di tutti i moduli messi insieme --- e implementa cinque
politiche, tre delle quali adattive.

Il capitolo procede cos\`i: prima il problema e la sua struttura analitica, poi
le cinque politiche in ordine di sofisticazione crescente, infine il ciclo che
le esegue. I risultati sperimentali che hanno motivato le scelte sono riassunti
qui dove servono a giustificare una decisione di progetto, e riportati per esteso
nel \cref{cap:25-valutazione-concorrenza}.

\section{Il problema}

\subsection{Che cosa fa un limite di concorrenza}

In \nf{}, il limite di concorrenza di una funzione \`e il valore contro cui
\code{FunctionQueueState.tryAcquireSlot()} confronta il contatore delle richieste
in volo (\cref{cap:09-async-queue}). Una richiesta che non ottiene lo slot resta
in coda.

Ne segue che il limite governa due grandezze contemporaneamente, e questa
duplicit\`a \`e la fonte di tutte le difficolt\`a che seguono.

\paragraph{Primo: l'utilizzazione.} Un limite troppo basso lascia le repliche
inattive mentre le richieste attendono. Un limite troppo alto le sovraccarica,
allungando il tempo di servizio di ciascuna richiesta.

\paragraph{Secondo: l'ammissione.} Una funzione pu\`o trattenere al pi\`u
$C + Q$ richieste, dove $C$ \`e il limite e $Q$ la dimensione della coda; oltre
quella soglia rifiuta. Il limite \`e quindi anche, insieme alla coda, la
finestra di ammissione.

La seconda propriet\`a \`e stata scoperta sperimentalmente, non progettata, ed \`e
costata un fallimento --- descritto nella \cref{sec:sojourn} --- in cui un
controllore ottimizzava correttamente la prima dimenticando la seconda, e
rifiutava il 4\% del traffico offerto.

\subsection{La struttura analitica}

Sia $\lambda$ il tasso di arrivo, $S(C)$ il tempo di servizio medio con limite
$C$, $W(C)$ il tempo di attesa in coda e $T(C) = W(C) + S(C)$ il tempo di
sojourn --- ci\`o che il chiamante osserva.

Dalla legge di Little~\cite{little1961proof} applicata al solo servizio, con il
sistema saturo, il throughput \`e $\mu(C) = C / S(C)$. Poich\'e le repliche
condividono CPU e memoria, $S$ cresce con $C$: pi\`u richieste contemporanee
significa pi\`u contesa, quindi

\[
  \frac{\mathrm{d}S}{\mathrm{d}C} > 0 \qquad \text{per ogni } C > 0 .
\]

Questa disuguaglianza ha una conseguenza immediata e distruttiva per una classe
intera di controllori:

\begin{quote}
\textbf{Un controllore che minimizza $S(C)$ non ha ottimo interno.} Poich\'e $S$
\`e monotona crescente, il gradiente punta in discesa ovunque, e una regola che
riduce il limite quando la latenza cresce cammina fino al proprio limite
inferiore.
\end{quote}

Il tempo di attesa ha il comportamento opposto: la coda drena pi\`u in fretta
quando pi\`u richieste sono servite in parallelo, quindi $W$ decresce con $C$. La
somma ha un minimo interno. La \cref{fig:sojourn} illustra la relazione.

\begin{figure}[htbp]
  \centering
  \resizebox{0.85\textwidth}{!}{% Schematic: why service time has no interior optimum and sojourn time does.
% Illustrative shapes, NOT measured data — the caption must say so.
\documentclass[border=8pt]{standalone}
\usepackage{nanofaas-figs}
\usepackage{pgfplots}
\pgfplotsset{compat=1.18}
\begin{document}
\begin{tikzpicture}
\begin{axis}[
    width=11cm, height=7cm,
    xlabel={limite di concorrenza $C$},
    ylabel={tempo (unit\`a arbitrarie)},
    xmin=1, xmax=10, ymin=0, ymax=13,
    domain=1:10, samples=200,
    axis lines=left, font=\sffamily\small,
    legend style={draw=nfslate!50, font=\sffamily\footnotesize, at={(0.98,0.98)}, anchor=north east},
    tick label style={font=\sffamily\footnotesize},
    grid=major, grid style={nfslate!20},
]
  % service time: monotonically increasing with concurrency
  \addplot[nfblue, line width=1.1pt] {2 + 0.55*x};
  \addlegendentry{$S(C)$ --- tempo di servizio}
  % wait: decreasing, ~ backlog / drain rate
  \addplot[nfamber, line width=1.1pt] {18/x};
  \addlegendentry{$W(C)$ --- tempo di attesa}
  % sojourn = sum, interior minimum
  \addplot[nfgreen, line width=1.5pt] {2 + 0.55*x + 18/x};
  \addlegendentry{$T(C)=W+S$ --- sojourn}
  % marker at the minimum: d/dC (0.55C + 18/C) = 0  ->  C = sqrt(18/0.55) ~ 5.72
  \addplot[only marks, mark=*, mark size=2.4pt, nfgreen]
      coordinates {(5.72,11.29)};
  \node[font=\sffamily\footnotesize, text=nfgreen, anchor=south]
      at (axis cs:5.72,11.5) {ginocchio};
  \draw[dashed, nfslate] (axis cs:5.72,0) -- (axis cs:5.72,11.29);
\end{axis}
\end{tikzpicture}
\end{document}
}
  \caption{Schema qualitativo (\emph{non} dati misurati). Il tempo di servizio
  cresce monotonamente con il limite di concorrenza e non ammette minimo interno;
  il tempo di attesa decresce; la loro somma --- il tempo di sojourn --- ha un
  ginocchio. Un controllore che decide da $S$ non ha nulla da trovare; uno che
  decide da $T$ s\`i.}
  \label{fig:sojourn}
\end{figure}

\subsection{Che cosa il control plane misura davvero}

La metrica \code{function\_latency\_ms} \`e misurata dal control plane fra
l'istante di dispatch e quello di completamento. \`E quindi il \emph{tempo di
servizio}, comprensivo del round trip HTTP, e non include l'attesa in coda.

La misura riportata nel \cref{cap:25-valutazione-concorrenza} quantifica lo
scarto: sotto il profilo a raffica il tempo di attesa medio \`e stato di
37--43\,ms contro un tempo di servizio di circa 5\,ms. La grandezza da cui i
primi controllori decidevano \`e circa \textbf{un settimo} di ci\`o che il
chiamante sperimenta.

La conseguenza \`e stata misurata e non solo temuta: in un confronto iniziale il
modo \code{BUDGETED} ha ridotto il tempo di servizio del 31\% su una funzione e
del 44\% sull'altra --- entrambe ampiamente dentro un SLO di 10\,ms --- e ha
\emph{peggiorato} il p95 end-to-end, perch\'e l'attesa creata dal suo stesso
limite ha pi\`u che assorbito il guadagno. Un limite di concorrenza non elimina
lavoro: lo sposta nel buffer.

\section{Le cinque politiche}

\begin{table}[htbp]
\centering
\caption{Le cinque politiche di controllo della concorrenza.}
\label{tab:modi}
\small
\begin{tabularx}{\textwidth}{@{}lXll@{}}
\toprule
\textbf{Modo} & \textbf{Regola} & \textbf{Segnale} & \textbf{Ambito} \\
\midrule
\code{FIXED} & Il limite configurato, invariato. & nessuno & per funzione \\
\code{STATIC\_PER\_POD} & repliche pronte $\times$ obiettivo per replica. &
repliche & per funzione \\
\code{ADAPTIVE\_PER\_POD} & Gradiente di latenza contro il minimo storico della
funzione. & servizio & per funzione \\
\code{BUDGETED} & Fabbisogno per l'SLO, servito da un budget di piattaforma con
equit\`a max-min pesata. & servizio, throughput & \textbf{globale} \\
\code{SOJOURN} & Modello: guardia di ammissione, legge di Little su tasso
predetto, orizzonte di drenaggio. & sojourn, servizio, coda & per funzione \\
\bottomrule
\end{tabularx}
\end{table}

Tutte le politiche sono soggette a due vincoli invalicabili, applicati a valle:
il limite effettivo non pu\`o superare quello configurato nella specifica
(\cref{cap:09-async-queue}), e non pu\`o scendere sotto il pavimento
configurato. La specifica della funzione resta il tetto e il pavimento; le
politiche si muovono all'interno.

\section{\code{FIXED} e \code{STATIC\_PER\_POD}}

\code{FIXED} \`e il default e non fa nulla: il limite \`e quello configurato. \`E
la configurazione di riferimento contro cui ogni altra viene confrontata.

\code{STATIC\_PER\_POD} \`e ventotto righe:

\begin{lstlisting}[language=Java]
int replicas = Math.max(1, readyReplicas);
int targetPerPod = control.targetInFlightPerPod() == null
        ? 1 : Math.max(1, control.targetInFlightPerPod());
long desired = (long) replicas * targetPerPod;
if (desired > configured) { return configured; }
return Math.max(1, (int) desired);
\end{lstlisting}

L'idea \`e che il limite sensato non \`e una propriet\`a della funzione ma della
capacit\`a disponibile: se ogni replica pu\`o servire bene due richieste
contemporanee e le repliche pronte sono cinque, il limite \`e dieci. \`E la
stessa nozione di \code{containerConcurrency} di Knative, con la differenza che
qui \`e il control plane a moltiplicare, avendo visibilit\`a sul numero di
repliche.

Il modo \`e utile e ha un limite noto: assume che l'obiettivo per replica sia
noto all'operatore. Determinarlo richiede un esperimento --- che \`e esattamente
ci\`o che le politiche adattive cercano di evitare.

\section{\code{ADAPTIVE\_PER\_POD}: il gradiente di latenza}

\subsection{Il segnale}

Il commento di classe enuncia la scelta del segnale con precisione:

\begin{quote}\itshape
Il segnale non \`e deliberatamente n\'e la profondit\`a di coda n\'e
l'utilizzazione. Una funzione che gira alla massima utilizzazione con una coda
vuota \`e sana, non satura, quindi nessuna delle due pu\`o localizzare il punto in
cui la concorrenza aggiuntiva inizia a costare throughput. Un tempo di servizio
crescente a lavoro costante \`e esattamente quel punto.
\end{quote}

L'osservazione \`e corretta e vale la pena isolarla, perch\'e distingue questa
famiglia di controllori da quelli basati su soglie di utilizzazione. Il segnale
\`e la \emph{degradazione} rispetto al miglior tempo di servizio osservato:

\[
  d = \frac{\bar{S}_{\text{int}} - S_{\text{base}}}{\bar{S}_{\text{int}}}
  \;\in\; [0, 1),
\]

dove $\bar{S}_{\text{int}}$ \`e la media dell'intervallo appena concluso e
$S_{\text{base}}$ il minimo storico. \`E la costruzione di TCP
Vegas~\cite{brakmo1995vegas} trasposta alle chiamate RPC, come nella libreria di
Netflix~\cite{netflixlimits}.

\subsection{La regola}

\begin{lstlisting}[language=Java]
if (degradation >= 0) {
    if (degradation >= highThreshold) {
        if (nowEpochMs - state.lastDecreaseEpochMs() >= downCooldown) {
            target = Math.max(minTarget, target - 1);
            state.lastDecreaseEpochMs(nowEpochMs);
        }
    } else if (degradation <= lowThreshold
            && nowEpochMs - state.lastIncreaseEpochMs() >= upCooldown) {
        target = Math.min(maxTarget, target + 1);
        state.lastIncreaseEpochMs(nowEpochMs);
    }
}
\end{lstlisting}

Passi unitari, con banda morta fra le due soglie (default 0{,}15 e 0{,}5) e
cooldown asimmetrici (30\,s in su, 60\,s in gi\`u). Il valore $-1$ della
degradazione codifica ``nessuna invocazione completata in questo intervallo'', e
in quel caso il limite non si muove: \emph{un intervallo vuoto non \`e evidenza di
niente}.

\subsection{La deriva della baseline}

\begin{lstlisting}[language=Java]
/**
 * ponytail: the baseline creeps up 0.2%/tick so a single lucky-fast interval cannot pin it
 * forever, while a real 2x degradation still stands out for minutes. Replace with a windowed
 * minimum if functions turn out to have long-period service-time seasonality.
 */
private static final double BASELINE_DRIFT_PER_TICK = 1.002;
\end{lstlisting}

Il minimo storico puro ha un difetto noto: un singolo intervallo fortunato ---
una finestra in cui il sistema era scarico --- fissa un riferimento che il
sistema non raggiunger\`a pi\`u, e da quel momento la degradazione appare
permanentemente alta. La deriva dello 0{,}2\% per tick fa dimenticare lentamente
i valori estremi. \`E, come il commento riconosce, la versione economica di un
minimo su finestra scorrevole.

\subsection{Perch\'e non basta}

Due difetti, entrambi misurati.

\paragraph{Non ha ottimo interno.} Il riferimento \`e il minimo osservato della
funzione stessa, e per la monotonia di $S(C)$ il gradiente punta in discesa
ovunque. Il commento di \code{SloDemandEstimator} lo registra: \emph{``a
controller that chases its own minimum has no interior optimum to find [\ldots]
the limit walks to its floor, which is what two runtimes were measured doing''}.

\paragraph{Decide su una risorsa condivisa senza vedere i vicini.} Con pi\`u
funzioni sullo stesso control plane, il limite di una si muove perch\'e un'altra
\`e arrivata, mentre il proprio carico non \`e cambiato. \`E il fenomeno misurato
nella campagna di co-tenancy (\cref{cap:25-valutazione-concorrenza}): la funzione
stabile perde il 28\% di throughput e guadagna il 36\% di latenza quando la
vicina compare. Il controllore interpreta il degrado come proprio e reagisce, ma
la causa \`e altrove.

I due difetti motivano rispettivamente i due modi successivi.

\section{\code{BUDGETED}: la concorrenza come risorsa di piattaforma}

\subsection{L'idea}

Il commento di classe la enuncia:

\begin{quote}\itshape
Ogni funzione dichiara ci\`o di cui ha bisogno, e una sola allocazione risponde a
tutte. [\ldots] Una decisione per funzione non pu\`o rispettare un budget
condiviso, perch\'e nessuna funzione sa che cosa stanno chiedendo le altre.
\end{quote}

La domanda viene quindi spezzata in due, e le due parti sono risolte da due
componenti distinti:

\begin{enumerate}
  \item \emph{Quanto serve a questa funzione} per rispettare il proprio SLO al
  carico che sta vedendo? Risponde \code{SloDemandEstimator}, per funzione.
  \item \emph{Quanto pu\`o dare la piattaforma}, dato che il budget \`e finito e
  altre funzioni chiedono? Risponde \code{ConcurrencyBudgetAllocator}, una volta
  per tick su tutte le richieste.
\end{enumerate}

Per costruzione la somma delle concessioni non pu\`o eccedere il budget, e le
funzioni non devono scoprirsi a vicenda attraverso l'interferenza.

\subsection{La stima del fabbisogno}

La regola base \`e il gradiente della libreria Netflix, con l'SLO come obiettivo
anzich\'e il minimo auto-misurato:

\[
  g = \operatorname{clamp}\!\left(\frac{L_{\text{target}}}{L_{\text{oss}}},\,
      0{,}5,\, 1\right),
  \qquad
  C_{\text{des}} = C \cdot g + \sqrt{C}.
\]

Tre elementi meritano commento.

\paragraph{L'obiettivo \`e un SLO, non un minimo osservato.} \`E la correzione
diretta del primo difetto di \code{ADAPTIVE\_PER\_POD}. Contro un obiettivo
fisso il gradiente \`e maggiore di 1 finch\'e la funzione \`e dentro il proprio
SLO, e minore di 1 solo quando la promessa \`e effettivamente violata. Esiste
quindi un punto di equilibrio interno.

\paragraph{Il pavimento a 0{,}5 sul gradiente.}
\begin{lstlisting}[language=Java]
/**
 * Floor on the gradient. Without it a single slow interval — a cold start, a GC pause — would
 * halve the limit repeatedly; with it the worst case is halving per tick, which recovers.
 */
private static final double MIN_GRADIENT = 0.5;
\end{lstlisting}
Limita il danno di un intervallo anomalo. Il caso peggiore diventa un
dimezzamento per tick, da cui il sistema si riprende, anzich\'e un crollo.

\paragraph{Il termine $\sqrt{C}$.} \`E margine deliberato. Un limite
dimensionato esattamente sul tasso di arrivo corrente non lascia nulla per una
raffica: il primo arrivo sopra la media si accoda e, se la coda \`e piena, viene
rifiutato. Per la legge di Little la concorrenza necessaria ad assorbire
$\lambda$ a tempo di servizio $S$ \`e $\lambda S$; il margine \`e ci\`o che sta
sopra, e cresce in modo sublineare cos\`i che un limite grande non porti uno
spreco proporzionalmente grande.

\subsection{Il tetto al ginocchio}

Il gradiente da solo non basta, e la ragione \`e misurata:

\begin{quote}\itshape
Un controllore guidato dall'SLO spende ogni millisecondo di latenza che gli viene
concesso: detto che pu\`o impiegare 10\,ms, \`e cresciuto fino a quattro volte la
concorrenza del modo a gradiente sullo stesso carico e ha consegnato il 6\% di
richieste in meno, perch\'e nulla nella regola chiedeva se la concorrenza
aggiuntiva comprasse dei completamenti.
\end{quote}

La correzione \`e un tetto alla crescita calcolato anzich\'e cercato:

\[
  C_{\text{ginocchio}} = \lambda_{\max} \cdot S_{\min}
  \;+\; \sqrt{\lambda_{\max} \cdot S_{\min}},
\]

ossia la legge di Little applicata ai due estremi che la funzione ha dimostrato
di saper raggiungere. Una volta che il throughput satura, $\lambda_{\max}$ smette
di crescere e il tetto smette di crescere con esso, qualunque cosa l'SLO
consentirebbe ancora.

Il commento identifica correttamente la parentela: \emph{``This is BBR's shape
rather than Vegas': track the two extremes separately and target their product,
instead of reading one degraded signal.''} Vegas legge un segnale degradato --- il
ritardo --- e ne deduce lo stato; qui si tracciano separatamente il miglior tempo
di servizio e il miglior throughput e se ne prende il prodotto.

Gli estremi decadono verso il presente all'1\% per tick, per la stessa ragione
per cui la baseline del modo precedente deriva: un minimo assoluto non si riprende
mai da un intervallo fortunato, e un massimo assoluto non ammette mai che la
macchina sia diventata pi\`u lenta.

Un'asimmetria importante \`e codificata esplicitamente:

\begin{lstlisting}[language=Java]
// Growth stops at the knee even when the SLO would allow more. Shrinking is never
// capped: an SLO breach is a broken promise, and no throughput argument outranks it.
if (desired > current) {
    desired = Math.min(desired, knee(spec.name(), observedLatencyMs, throughputRps));
}
\end{lstlisting}

Il tetto vincola solo la crescita. La contrazione non \`e mai limitata, perch\'e
una violazione dell'SLO \`e una promessa infranta e nessun argomento di
throughput la sopravanza. \`E una gerarchia di obiettivi dichiarata nel codice,
non implicita.

\subsection{L'allocazione: equit\`a max-min pesata}

L'allocatore implementa l'equit\`a max-min pesata~\cite{bertsekas1992datanetworks}:

\begin{enumerate}[leftmargin=*]
  \item Assegna a ciascuna funzione il proprio pavimento.
  \item Distribuisce il residuo proporzionalmente ai pesi, ma nessuna funzione
  riceve pi\`u di quanto ha chiesto.
  \item Ripete finch\'e resta budget e qualcuno \`e ancora insoddisfatto.
  \item Se le quote arrotondano a zero, distribuisce il residuo a unit\`a singole
  anzich\'e ciclare a vuoto.
\end{enumerate}

Il commento giustifica la scelta rispetto all'alternativa ovvia --- scalare ogni
richiesta dello stesso fattore --- con due propriet\`a:

\begin{quote}\itshape
Una funzione che chiede meno della propria quota non viene mai tagliata, quindi
una funzione piccola non \`e punita per aver condiviso la piattaforma con una
grande. E nulla resta sul tavolo mentre qualcuno \`e ancora a corto: lo
scalamento proporzionale lascerebbe inattiva capacit\`a che una funzione affamata
potrebbe usare.
\end{quote}

\subsection{I pavimenti prima di tutto}

Il trattamento del caso in cui i soli pavimenti eccedono il budget \`e la
decisione pi\`u interessante dell'allocatore:

\begin{quote}\itshape
Una funzione a cui viene concesso zero non pu\`o servire, e una piattaforma che
affama una funzione fino al silenzio totale per soddisfarne una vicina ha preso
una decisione di ammissione fingendo di prenderne una di scheduling. Se i soli
pavimenti eccedono il budget, il budget \`e per definizione troppo piccolo per le
funzioni registrate, e ogni pavimento viene concesso: l'alternativa \`e rompere
una promessa che la piattaforma fa a ogni tenant per tenere in ordine
l'aritmetica.
\end{quote}

Il budget viene quindi sforato. \`E la scelta corretta: un budget \`e uno
strumento di ripartizione, non un vincolo fisico, e quando diventa
insoddisfacibile il fatto da segnalare \`e che \`e mal dimensionato, non che una
funzione debba smettere di esistere.

\subsection{Il budget di piattaforma}

\begin{lstlisting}[language=Java]
/**
 * A platform capacity statement, not something to infer from load: derived from observed
 * throughput it would grow under load and shrink when idle, which is the opposite of what a
 * budget is for. Unset, it scales with the cores the control plane can see, which is a guess —
 * but a guess that at least tracks the machine rather than a constant.
 */
public int totalBudgetOrDefault() {
    if (totalBudget != null && totalBudget > 0) { return totalBudget; }
    return Math.max(8, Runtime.getRuntime().availableProcessors() * 4);
}
\end{lstlisting}

\`E lo stesso principio dell'SLO di default (\cref{cap:06-nucleo}): una grandezza
che esprime un vincolo non deve essere derivata dall'osservazione, o cesser\`a di
vincolare esattamente quando serve.

\section{\code{SOJOURN}: decidere dalla latenza del chiamante}
\label{sec:sojourn}

Il modo \code{SOJOURN} \`e il pi\`u recente e il pi\`u istruttivo, perch\'e la
sua documentazione conserva due progetti falliti insieme a quello adottato.

\subsection{Primo fallimento: il gradiente sulla latenza end-to-end}

L'idea immediata --- alimentare la regola a gradiente esistente con la latenza
end-to-end anzich\'e con il tempo di servizio --- \`e instabile, e \`e stata
respinta prima dell'implementazione. La regola riduce il limite quando la latenza
eccede l'obiettivo; un limite pi\`u piccolo drena la coda pi\`u lentamente, quindi
l'attesa cresce, quindi il limite si riduce ancora, fino al pavimento. \`E un
anello di retroazione positiva.

\subsection{Secondo fallimento: la ricerca del minimo}

Poich\'e $T(C)$ ha un minimo interno (\cref{fig:sojourn}), cercarlo per passi
successivi --- muovi, confronta con l'ultima volta, prosegui se \`e migliorato,
inverti se \`e peggiorato --- \`e ben posto in linea di principio. \`E stato
implementato e falsificato dalla misura:

\begin{center}
\small
\begin{adjustbox}{max width=\textwidth}
\begin{tabular}{@{}lrr@{}}
\toprule
 & \textbf{\code{ADAPTIVE\_PER\_POD}} & \textbf{\code{SOJOURN}, a ricerca} \\
\midrule
richieste rifiutate & 32 (0{,}004\%) & \textbf{36\,181 (4{,}275\%)} \\
limite medio [min--max] & 7{,}8 [4--8] & 4{,}0 [1--8] \\
\bottomrule
\end{tabular}
\end{adjustbox}
\end{center}

La causa \`e un problema di attribuzione:

\begin{quote}\itshape
Non pu\`o distinguere ``la mia mossa ha aiutato'' da ``il carico \`e calato''. In
un avvallamento ha attribuito il sollievo naturale alla propria decisione di
restringere, ed era ancora piccolo quando il carico \`e tornato. Confrontare due
istanti presuppone che il mondo sia rimasto fermo fra di essi.
\end{quote}

Il dettaglio pi\`u interessante \`e che il punto di lavoro scelto non era
sbagliato: a limite 4 la ricerca consegnava il 96{,}5\% del throughput con met\`a
del tempo di servizio. Ci\`o che ha mancato \`e la seconda funzione del limite ---
essere la finestra di ammissione --- e dimensionarlo dalla sola latenza ha fatto
perdere il 4\% del traffico.

\subsection{Il progetto adottato: calcolare anzich\'e confrontare}

Il controllore adottato non confronta due intervalli. Calcola, in tre parti con
un ordine di autorit\`a stretto.

\paragraph{1. La guardia di ammissione.} Ha la precedenza assoluta:

\begin{lstlisting}[language=Java]
private static final double HIGH_WATER = 0.7;

if (underQueuePressure(queueDepth, bounds)) {
    // Never below what is already in flight while the buffer is filling: lowering the limit
    // here shrinks the admission window and refuses the very traffic that is queueing.
    needed = Math.max(needed, Math.max(inFlight + 1, bounds.floor()));
}
\end{lstlisting}

Quando la coda supera il 70\% della propria capacit\`a, il limite non pu\`o
scendere, qualunque cosa l'aritmetica della latenza preferisca. Il commento \`e
esplicito sul fatto che questa non \`e una preferenza tarabile: \emph{``refusing
traffic is worse than being slow, and that ordering is not a tuning
parameter''}.

\paragraph{2. La legge di Little su un tasso predetto.}

\[
  C_{\text{servizio}} = \hat{\lambda} \cdot S,
  \qquad
  C_{\text{drenaggio}} = \frac{Q \cdot S}{H},
  \qquad
  C = \left\lceil C_{\text{servizio}} + C_{\text{drenaggio}} \right\rceil
\]

dove $\hat{\lambda}$ \`e il tasso di arrivo \emph{predetto}, $Q$ la profondit\`a
di coda corrente e $H$ l'orizzonte di drenaggio concesso.

\paragraph{3. La promessa come orizzonte.} Quando il sojourn osservato eccede
l'obiettivo, l'orizzonte di drenaggio si accorcia in proporzione:

\[
  H = \max\left(1\,\mathrm{s},\; 5\,\mathrm{s} \cdot
      \frac{T_{\text{target}}}{T_{\text{oss}}}\right).
\]

Un orizzonte pi\`u corto richiede pi\`u concorrenza per smaltire lo stesso
arretrato. La logica \`e che l'attesa \`e l'unica parte della latenza del
chiamante che un limite di concorrenza possa effettivamente accorciare.

\subsection{La predizione del carico}

La predizione $\hat{\lambda}$ \`e ci\`o che elimina il problema di attribuzione:
il controllore smette di confrontare istanti e inizia a calcolare da un modello.
\code{LoadForecast} implementa il metodo lineare di Holt con trend smorzato:

\[
  \ell_t = \alpha \lambda_t + (1-\alpha)(\ell_{t-1} + \phi b_{t-1}),
  \quad
  b_t = \beta(\ell_t - \ell_{t-1}) + (1-\beta)\phi b_{t-1},
  \quad
  \hat{\lambda}_{t+1} = \ell_t + \phi b_t,
\]

con $\alpha = 0{,}5$, $\beta = 0{,}3$, $\phi = 0{,}8$.

La scelta \`e giustificata cos\`i:

\begin{quote}\itshape
Una media mobile \`e in ritardo su una rampa, e una rampa \`e dove la coda si
costruisce: quando una stima in ritardo se ne accorge, il buffer \`e gi\`a pieno.
Il termine di trend vede la pendenza; lo smorzamento impedisce di portare quella
pendenza oltre la cima della rampa, dove un'estrapolazione non smorzata
sovrastima di pi\`u.
\end{quote}

Il primo campione produce livello senza trend: \emph{``una lettura \`e un livello
e nessun trend; una pendenza richiede due punti, e inventarne uno farebbe
predire una rampa al primo intervallo di ogni funzione''}.

\subsection{Il risultato}

Sotto arrivi aperti --- richieste programmate dall'orologio anzich\'e da un
generatore a ciclo chiuso --- il confronto con \code{ADAPTIVE\_PER\_POD} sullo
stesso carico:

\begin{center}
\small
\begin{adjustbox}{max width=\textwidth}
\begin{tabular}{@{}lrrr@{}}
\toprule
 & \textbf{\code{ADAPTIVE\_PER\_POD}} & \textbf{\code{SOJOURN}} & \\
\midrule
servite & 549\,856 & \textbf{567\,210} & $+3{,}2\%$ \\
rifiutate & 110\,117 (16{,}7\%) & \textbf{92\,785 (14{,}1\%)} & $-15{,}7\%$ \\
p95 end-to-end & 120{,}48\,ms & \textbf{107{,}68\,ms} & $-10{,}6\%$ \\
p99 end-to-end & 157{,}67\,ms & \textbf{133{,}82\,ms} & $-15{,}1\%$ \\
attesa in coda & 35{,}61\,ms & 31{,}03\,ms & $-12{,}9\%$ \\
tempo di servizio & 1{,}69\,ms & 2{,}00\,ms & $+18\%$ \\
limite medio & 3{,}9 & 2{,}8 & \\
\bottomrule
\end{tabular}
\end{adjustbox}
\end{center}

Lo scambio \`e andato nella direzione prevista: il tempo di servizio peggiora del
18\%, l'attesa migliora del 13\%, e poich\'e l'attesa \`e venti volte il tempo di
servizio il chiamante ci guadagna. Il dettaglio istruttivo \`e che \code{SOJOURN}
ottiene questo con \emph{meno} concorrenza media: non vince aggiungendo
capacit\`a ma collocandola nel tempo, che \`e ci\`o per cui la predizione era
stata introdotta.

Il \cref{cap:25-valutazione-concorrenza} riporta i caveat --- una sola esecuzione
per modo, entrambi i sistemi in sovraccarico --- e il
\cref{cap:27-discussione} li tratta come limiti alla validit\`a.

\section{Il governor: il ciclo che esegue le politiche}

\code{ConcurrencyGovernor} \`e il componente che a ogni tick (default 5\,s)
raccoglie le osservazioni, invoca la politica pertinente per ciascuna funzione, e
pubblica i limiti attraverso \code{ScalingMetricsSource}.

Ha una struttura obbligata dalla differenza di ambito fra le politiche. I modi
\code{STATIC\_PER\_POD} e \code{ADAPTIVE\_PER\_POD} decidono per funzione e sono
gestiti da \code{ConcurrencyControlCoordinator} in un ciclo. Il modo
\code{BUDGETED} decide su tutte le funzioni insieme e riceve quindi la lista
completa delle osservazioni in una sola chiamata. Il modo \code{SOJOURN} torna a
decidere per funzione, ma legge tre segnali anzich\'e uno.

\`E il motivo per cui l'interfaccia dei controllori non \`e uniforme: un modo che
alloca una risorsa condivisa non pu\`o essere ricondotto alla stessa firma di uno
che decide in isolamento senza inventare uno stato globale nascosto.

\subsection{L'intervallo come unit\`a di osservazione}

Tutti i controllori adattivi leggono i timer cumulativi di Micrometer e li
convertono in medie dell'intervallo appena concluso, attraverso
\code{LatencyIntervals}:

\begin{lstlisting}[language=Java]
/**
 * Turns a cumulative service-time timer into the mean of the interval just ended.
 *
 * <p>The cumulative value answers "since startup", which flattens exactly the change a controller
 * reacts to: a function that has served a million fast requests and is now slow still reads fast.
 * Only the difference between two readings describes the present.</p>
 */
\end{lstlisting}

\`E lo stesso errore, e la stessa correzione, gi\`a incontrati per la metrica
\code{rps} dell'autoscaler (\cref{cap:11-autoscaler}). Che il progetto abbia
dovuto risolverlo due volte in due moduli \`e un piccolo indizio a favore di una
astrazione condivisa che non esiste.

\section{Configurazione}

\begin{center}
\small
\begin{adjustbox}{max width=\textwidth}
\begin{tabular}{@{}llp{6cm}@{}}
\toprule
\textbf{Chiave} & \textbf{Default} & \textbf{Significato} \\
\midrule
\code{nanofaas.concurrency-control.poll-interval-ms} & 5\,000 & Periodo del tick. \\
\code{...default-target-in-flight-per-pod} & 2 & Obiettivo iniziale. \\
\code{...total-budget} & $\max(8, 4 \times \text{core})$ & Budget globale per
\code{BUDGETED}. \\
\midrule
\multicolumn{3}{@{}l@{}}{\emph{Per funzione, in \code{scalingConfig.concurrencyControl}:}} \\
\code{mode} & \code{FIXED} & \\
\code{targetInFlightPerPod} & 2 & \code{STATIC}/\code{ADAPTIVE}. \\
\code{minTargetInFlightPerPod} & 1 & Pavimento. \\
\code{maxTargetInFlightPerPod} & 8 & Tetto. \\
\code{upscaleCooldownMs} & 30\,000 & \code{ADAPTIVE}. \\
\code{downscaleCooldownMs} & 60\,000 & \code{ADAPTIVE}. \\
\code{highLoadThreshold} & 0{,}5 & \code{ADAPTIVE}. \\
\code{lowLoadThreshold} & 0{,}15 & \code{ADAPTIVE}. \\
\code{targetLatencyMs} & 250 & SLO: di servizio per \code{BUDGETED},
end-to-end per \code{SOJOURN}. \\
\code{weight} & 1{,}0 & Peso nell'allocazione \code{BUDGETED}. \\
\bottomrule
\end{tabular}
\end{adjustbox}
\end{center}

Va notato che \code{targetLatencyMs} cambia significato fra i due modi che lo
usano: in \code{BUDGETED} \`e un SLO sul tempo di servizio, in \code{SOJOURN}
sull'end-to-end. \`E un'ambiguit\`a reale nel modello di configurazione, mitigata
dal fatto che ciascun modo interpreta il campo coerentemente con la propria
semantica dichiarata, e non risolta.

\section{Sintesi}

\begin{table}[htbp]
\centering
\caption{Le cinque politiche: che problema risolve ciascuna, e quale lascia
aperto.}
\label{tab:modi-sintesi}
\small
\begin{tabularx}{\textwidth}{@{}lXX@{}}
\toprule
\textbf{Modo} & \textbf{Risolve} & \textbf{Lascia aperto} \\
\midrule
\code{FIXED} & Nulla; \`e il riferimento. & Tutto. \\
\code{STATIC\_PER\_POD} & Il limite segue la capacit\`a disponibile. &
L'obiettivo per replica va determinato sperimentalmente. \\
\code{ADAPTIVE\_PER\_POD} & Trova il proprio punto senza taratura. & Non ha
ottimo interno; ignora i vicini. \\
\code{BUDGETED} & La risorsa condivisa \`e allocata una volta sola, con
equit\`a. & Decide ancora dal tempo di servizio. \\
\code{SOJOURN} & Decide da ci\`o che il chiamante osserva; non confronta
istanti. & Mai osservato fuori dal sovraccarico; guardia al 70\% non tarata. \\
\bottomrule
\end{tabularx}
\end{table}

Le due colonne di destra sono la struttura del capitolo: ogni modo esiste perch\'e
il precedente ha lasciato aperto qualcosa, e ciascuno lascia aperto qualcosa a
sua volta. La progressione non \`e terminata --- l'ultima riga della tabella
elenca due questioni che il \cref{cap:28-lavori-futuri} riprende.
